\chapter{Story}

\begin{multicols}{2}

This chapter contains skills and tools that are designed to help the players create a story together. 

\textbf{Contacts and reputation:} The contacts and leverages systems are meant to track story progress in terms of social progression.

\textbf{Playing without Leverages and Contacts} Simply consider that every NPC from the contacts chapter offers their services at base rate. Conduct the social part of the game narratively only with social tests. Use missions without any of the numerical social rewards and just reward karma and money.

\section{Knowledge}

\skill{Knowledge}{ 2x Knowledge or 2x INT}{

\textbf{Knowledge test:} is equivalent to asking your character if they know or remember something from their studies. Use 2x the relevant knowledge or 2x INT as the skill value. This test cannot be done twice for the same thing.

The player is welcome to suggest what the character is supposed to know from the test. The GM takes that information into consideration when choosing the test DL. If the test is a hit or crit, what the player said is true, or partially true.

If the player suggests nothing and just asks if the character knows something, consider the topic and how likely is any information about it supposed to be known to choose a DL. If the test result was a crit or hit, come up with a piece of information. If it was a graze or miss, tell them nothing is known or give them misinformation.

\textbf{Deduction test:} If the necessary facts for deduction are already known, you can use 2x INT as the skill value. Whenever a player claims that the character knows something, the GM can ask for a knowledge test to determine if the character really knows it.

To use this skill, make a knowledge test against an arbitrary difficulty, representing the challenge of reaching a good conclusion from the available information. Critical successes mean that the character can make the best deduction from the information they have, even determining what information is missing to make a better deduction and remembering everything they know. Successes mean they have a good guess and remember most of what they know. Grazes mean that you don't know or can't remember, but you might know where to look. Misses mean that you reach an incorrect conclusion or a different one from the correct one, or you can't remember anything.
}

\subsection{Specific Knowledge}

Specific knowledge is about things that pertain the setting or adventure. It is knowledge about languages, people, places, cultures, and whatever else can be thought. 

\textbf{Why have this?}

\begin{proplist}
    \item \textbf{Substitute for area of knowledge:} can be used instead of an area of knowledge when it is higher and applies to the situation. Can also be used standalone as a knowledge when no areas are applicable.
    \item \textbf{Requirement/Enablement for specific actions:} some actions are vetoed or get heavy penalties for not matching a knowledge level.
    \item \textbf{Memory of character experiences:} stays in the sheet as a memento of the character's experiences. Enables creative usage in the future.
\end{proplist}

\textbf{knowledge levels:} specific knowledge has levels, just like areas of knowledge. A knowledge test can increase the knowledge level by a certain amount up to a specific limit. Shallow learning experiences can only increase knowledge up to a low level, like 1 or 2. More intensive and prolonged ones increase it further and by greater amounts.

\textbf{Gaining specific knowledge:} happens through a knowledge or skill test. The test has a DL and the result describes how much knowledge is gained, from near nothing on a miss to a lot on a crit.

\textbf{Learning tests:} Can be downtime activities, like spending a month in a town socializing. Can be exploration events, like climbing a tall tree to see the surrounding area. Can be a social interaction, like when someone explains how to behave in a social situation. The better the learning experience, the lower the DL and/ or the higher the potential level gain. The DL and learning potential for this are chosen subjectively.

\textbf{Example learning tests:}

\begin{proplist}
    \item spending a month in a town socializing has a persuasion or insight DL of 5 and can increase local knowledge by 1/2/3/4 up to a maximum of level 4.
    \item climbing a tall tree to see the surrounding area has a geography DL of 3 and can increase local knowledge by 0/1/2/2 up to a maximum of level 2.
    \item spending a month learning a new language is a linguistics test with a DL of 10 and can increase language knowledge by 0/1/2/3 up to a maximum of level 4.
\end{proplist}

Examples of varied knowledge:

\textbf{Location - place:} Knowledge about the people in the town and their local customs. \\
\textbf{Survival - region:} Knowledge about a specific region, its roads and paths, food sources, water sources, natural disasters, economic resources, geographical landmarks, etc. This type of knowledge is often used alongside survival actions like traveling and searching. \\
\textbf{Language:} Knowledge about a specific language. Having level 3 in a language is necessary for fluent speech. \\
\textbf{Art - type:} Knowledge about a specific form of art, for example, music, its greatest artists, and works. \\
\textbf{Faction - name:} Knowledge about the people of a faction, their history, politics, wealth, alliances, etiquette, etc. \\
\textbf{Religion - name:} Knowledge about the gods of a religion and their domains. \\
\textbf{Animal handling - animal:} Knowledge about taming and riding a specific animal or type of animal. Some animals may not be tameable without this. \\
\textbf{Driving - vehicle:} May be necessary to be able to drive a special vehicle. \\

\section{Social Scenes:}

Social scenes in this game are just communication between characters. The objective of the system is to add depth to the exchange, not to promote a well defined game. 

The social test offers a way to structure the scenes so that they become collaborative storytelling.

\subsection{Social tests}
Social tests are skill tests between persuasion, deception, will, insight, cunning and knowledge. Which skill is used to attack and defend is determined by the context of the action and the GM's decision. 

\textbf{When to call for one:} Social tests are called to decide the course of a relevant change in the scene. Examples of relevant changes are: someone now believes something, a trade was settled at a given value, putting a suspicious NPC at ease, taunting a guard to attack you and succeeding. The point is that for a test to be called, there must be a consequence. 

A test is always a way to add uncertainty to an event, that is why it has dice. It is used whenever the GM does not want to say outright what the outcome of something was. This is important, because it gives players a chance to add to the story.

The character's skills are not the only variable in play in social interactions. Someone could make a persuasion test to attempt to get someone else to sell them their dagger, but the same test can be used to convince them of murdering their children. Both tests would be Persuasion vs Insight, but the odds should be vastly different. That is why an Interpretation bonus is introduced. 

\textbf{Interpretation Bonus (IB):} This is a value added to a social action test that represents the chances of it working out. The bonuses and penalties should range between -5 to +5. Larger values can be used instead of a straight veto or auto pass.

The GM should be careful giving positive IB for acting performance, because the character's skills of persuasion and deception already capture that component and should not be overriden. A performance may still help increasing the IB to the extent that the details convince the GM that the action is more likely to succeed than they otherwise would think.  

\textbf{Choosing an IB:} 

\begin{itemize}
  \item An IB of 10 is almost a guaranteed success, it should be given only in situations where the NPC should want to agree with the player anyways. Often, it is better to just allow the player to skip the test and get the desired outcome, but if the GM wants to keep the test, this is the IB to use.
  \item IBs around +5 are the most common ones for situations where a character is very likely to agree unless there is some special reason not to. A failure in such a test is just a prompt to find such a special condition to failure. Use this when player has made a very good argument and good acting. They are a reward for good roleplay, reading the scene well and coming up with practical game solutions.
  \item IBs between +2 and -2 are for common situations, usually improvised, where the one character does not have a strong opinion on the matter. An example would be to ask for a free pint of beer in a tavern. Give a bonus or penalty to based on how invested the player is in the interaction and how they show that in roleplay.
  \item IBs around -5 are for situations where the player is trying to do something that the NPC does not want to do, but it is not completely out of the question. For example, asking a shopkeeper for a discount out of the blue may warrant an IB of -5, but giving them a descent reason to do it may reduce the penalty to -3 or lower. This is a penalty to someone forcing the action where it does not belong or when the player has said something really bad in the roleplay, like insulting the NPC or being completely out of character. A success in such an action is a prompt to come up with a good reason why that should work.
  \item Finally, -10 is just to avoid giving straight vetos. It's an answer to at least let them try to do it.
\end{itemize}

Two ways of performing this test, one focusing on character skill, the other on player skill:

\textbf{Attacker rolls:} the attacker adds their skill value, IB and dice. The defender tells their defense skill value and must act according to the result in the dice. The dice result is public, so everyone knows how well the test went and what the result is. This is the simpler version, used for quickness and clarity. It abstracts the scene and gives more importance to character skill.

\textbf{Defender rolls:} The attacker tells the defender their skill value and the IB, which becomes the DL. The defender rolls with their respective skill against it. The degree of success does not need to be disclosed to the other players. This means that the attacker will have to interpret the defender's reaction to determine how well they did. When the defender hits or crits, they can also choose to give a false impression of their degree of success, like pretending to be afraid, or to believe in something without the attacker knowing. This type of roll requires more effort in acting. Use it to emphasize player skill.

\textbf{Test Outcomes:} The outcome of a social action depends on what skill is used for the test and the context of the action. An intimidation that goes wrong will have worse consequences than an attempt to seduce someone by giving a gift. The consequences should be aligned with the degree of success or failure. Maybe a critical success includes a bonus, or maybe a critical failure makes the NPC hostile. The GM is the ultimate adjudicator of the consequences.

\textbf{Repeating attempts:} the rule that a test should change the scene in some way already reinforces that a test to achieve the same objective should not be repeated. Sometimes it is clear that the interaction is a success or is ruined. Other times, the change is more subtle: you try convincing someone through logic and they don't agree, then you throw in a lie to convince them. The idea for repeating social tests is that each next test to achieve the same objective must increase the stakes, until the result forces a chance in the scene.

\textbf{Social interactions in combat:} If the social action is performed within a combat context, it should be used within the rules of the Morale system.

\subsection{Common uses}

These are guidelines on how to use the social skills. They are informative, and exist to illustrate usage of the rules rather than determine it.  

\textbf{Negotiating, Convincing, and Explaining - persuasion or deception against insight or knowledge:} These are the most common social actions. If they want to focus on the vibes to get a sense of the other person's intentions, they will use insight. If they want to check the logic of the argument, they will use knowledge. The description of what is discovered from a successful defense should be different for each defense.  

\textbf{Negotiating prices:} When negotiating, a crit gives contact prices to non contacts and 25\% discount to contacts, a hit gives regular prices, a graze gives the increased price to non contacts and normal price to contacts and a miss refuses negotiation with non contacts and loses 1 trust with contacts. 

\textbf{Lying and falsehoods - insight vs deception:} When anyone gives an information, an insight vs deception check can be called to determine whether the information is believed or not. Anyone can call for that test, but when the GM does, it must be done. On a crit or hit, the one who used insight knows the other one's true intentions, but not necessarily the truth of what they said. On a miss or graze, the one who used insight is convinced that the other one is telling the truth, even if they are lying. 

\textbf{Persuasion vs Deception:} Sometimes it can be ambiguous whether a social action is persuasion or deception. As a rule of thumb, the more important the lie is for the success of an interaction, the greater the indication that deception should be used. On the other hand, if the focus is in the overall story, feelings, or the relationship between the characters, persuasion is more appropriate. For example, if you are trying to convince someone to give you money by telling them that you need it to pay for your sick mother's medicine, but none of that is true, that is a deception. If it were true, it would be persuasion. But maybe they don't care about your mother, and they will just say that they are sorry and move on, even in a critical success. In that case, you have to pass two tests, but use only one dice roll for both. The first test is to see if you can convince them that you are telling the truth, and the second test is to see if they care about what you are saying.  

\textbf{Changing emotional state:} Social actions can be used to make a character afraid, enraged, or distracted. The GM should determine the appropriate skill to use for each case. 

\textbf{Intimidation - will, persuasion vs will, insight, knowledge:} The usual intimidation attempt is meant to make someone afraid and is made as a will vs will test. For example, when you look someone in the eye with murderous intent, you can use will vs will to make them afraid. Explaining all the ways someone's business is going to fail and telling them just how they are going to be destitute is an intimidation that uses persuasion against will, insight or even knowledge, depending on whether the defender wants to rely on their will to fight, their intuition that the other is bluffing or their knowledge of the situation. The IB should reflect the physical reality of the situation. A giant armed muscled and armored barbarian is obviously scarier than a 10 year old kid. But a dragon probably finds both quite unthreatening. This difference should be observed in the IB.

\textbf{Provocation - persuasion or deception vs insight:} When you insult someone, you can use persuasion vs insight or will to make them enraged. As a recommendation, insults are most effective when they exploit someone's frustrations or pride. Subtlety and timing are important. 

\textbf{Seduction - persuasion vs insight:} persuasion vs insight or will is used to seduce someone. This can be used as a tool to get things characters would otherwise not have access to, like being invited places and receiving gifts. 

\textbf{Distraction - persuasion vs cunning, insight:} When distracting someone, you can use persuasion vs cunning to make them focus attention on you. It can be used in a situation where you are just making a show, telling someone a story, starting a discussion or doing any long winded kind of interaction to capture attention. 

\textbf{Disguise - deception vs insight or cunning:} Use knowledge of a faction, culture or language as enabler with a deception test when trying to pretend being one of them. Disguise only covers the behavioral part. Dressing up and owning specific items is not accounted for.

\textbf{Interrogate - insight vs deception:} Ask a question, then roll Insight vs Deception to find out if someone is telling the truth. IB depends on the question and interrogation conditions. 

\textbf{Group tests:} are made with the skill of a single character. Other characters help by being a part of the scene and modifying the IB as a consequence, but the bonus is subjective.

\textbf{Using skills to skip tedious social procedures:} Sometimes players may not want to go through the roleplay of a social scene. In these cases, they can use usually persuasion but sometimes other appropriate skill to represent their social skill and skip the scene and go to the outcome. The player should give a general description of what they are trying to achieve, and the GM should determine the IB based on how well it makes sense. A hit or crit means that the desired outcome is achieved.


\section{Leverages}

Leverages are social levers used to gain access to goods, services, missions, jobs and any other social advantages.


\subsection{Renown} 
Ranges from 0 to 10 \\
Renown means that the character is considered more competent and important. It measures what is the equivalent of a social status of a character. \\
Characters with high renown will be better paid, will be called for more important tasks, and will have access to more exclusive areas of society. \\

\begin{stattable}{Renown scale}{@{}lc@{}}
Renown & Recognition  \\
\midrule
\rowcolor{fallowtint}
+10 & Legendary \\[3pt]
+8 to +9 & Leader \\[3pt]
+5 to +7 & Noble \\[3pt]
+2 to +4 & Skilled \\[3pt]
0 to +1 & Nobody \\[3pt]
\end{stattable}

\textbf{Intimidation:} The higher the renown, the higher the social status of an NPC that can be intimidated without forceful coercion. 

\subsection{Honor}  
Ranges from -10 to 10. \\
Low Honor is the reputation for being cruel, aggressive, or a liar.  \\
High Honor is the reputation for altruistic, good and well meaning.  \\
Characters with low Honor will be called for immoral or illegal tasks, while those with high Honor will be called for tasks to help people, receive gift and decorations.\\

\begin{stattable}{Honor scale}{@{}lcc@{}}
Honor & Recognition & Consequence \\
\midrule
\rowcolor{fallowtint}
+10 & Divine & +3 persuasion \\[3pt]
+7 to +9 & Hero & +1 persuasion \\[3pt]
+3 to +6 & Goodman & treated as contact overall\\[3pt]
-2 to +2 & Normal & Nothing\\[3pt]
-3 to -6 & Feared & Intimidating to lower social status \\[3pt]
-7 to -9 & Lawless & Law does not protect, Non-contacts reject \\[3pt]
-10 & Monster & actively hunted \\[3pt]
\end{stattable}

Effects are cumulative from 0 upwards and downwards. \\ 


\subsection{Popularity} is tracked from city to city or region to region and it defines whether the character's renown and honor is recognized there. The popularity scale is not necessary for someone to know a character. Some characters are more well informed than others and can know things for other reasons. Popularity sets the general reaction of the public.
The base amount of popularity is equal to the strength of the most influential Influencer that is a contact. Celebrating and living lavishly can override this base, if higher.

\begin{stattable}{Popularity scale}{@{}lc@{}}
Popularity & Recognition  \\
\midrule
\rowcolor{fallowtint}
+8 to +10 & All known \\[3pt]
+5 to +7 & Renown/Honor > 3 or < -3 are known \\[3pt]
+2 to +4 & Renown/Honor > 7 or < -7 are known \\[3pt]
0 to +1 & Nobody knows \\[3pt]
\end{stattable}




\subsection{Specific leverages}

\textbf{Affiliation:} Having/Not having proven allegiance with a contact. Coalitions and alliances exist. People will like whoever is on their side and dislike who isn't.

\begin{proplist}
    \item \textbf{Allegiance:} Affiliation can be used as a choice between two factions, where allying with one caps trust with the other.
    \item \textbf{Representation:} proving to represent by someone allows using their renown and Honor instead of the character's.  
\end{proplist}

\textbf{Titles:} Titles are in-setting tokens, like title of knighthood, insignias of a guild, knightly armor, a branded sword and other tangible proofs of status and affiliation. With those, anyone that recognizes the title automatically recognizes the character as one of a given Renown. There is no particular rule for those, since they are setting specific. 


\subsection{Tags}

Some NPCs have tags that define how they react in the game. It is possible to find out an NPCs tags with an insight test vs deception. 

\begin{proplist}
    \item \textbf{Seduceable:} If character is open for romantic relationship. Trying to seduce one that isn't receives automatic -5 seduction.
    \item \textbf{Etiquette:} groups require knowledge of their customs to allow access. Requires disguise to blend in.
    \item \textbf{Moralist:} Treats good characters with trust +1 or bad ones with -1 per Honor recognition level before being a contact.
\end{proplist}

\subsection{Wealth and Money}

Currency is like real world currency. It is used for trading and counted in coins.

\textbf{Currency:} Copper, silver, and gold are everyday currencies. Ten coppers are worth the same as one silver, and 10 silvers are worth the same as 1 gold. 1 coin weighs 10g, which means 50 coins weigh 500g, equivalent to 1 small item.

\textbf{Assets:} These are things one can claim to own. Assets represent land, rare items, slaves, investments, a business, or another form of wealth that is owned but has a fluctuating value compared to regular currency.


\section{Contacts}

Contacts provide information, utilities, and services. Those who are not contacts usually charge for exclusive information, charge higher prices for the same goods, hide special stock, refuse to perform certain tasks, etc.\\

\textbf{Acquiring Contacts:} A contact is any NPC with whom a character has had trust once. Old and unimportant contacts can be scratched from the sheet.

Contacts have a measure of relationship: Trust. They vary from -3 to +3. 

\textbf{Gaining trust:} This can be done by:

\begin{proplist}
    \item \textbf{completing missions:} Completing a mission for someone increases trust with them.
    \item \textbf{succeeding at work with/for someone:} Make a test of the most relevant skill or knowledge for the work against a DL to gain trust. On a crit, gain 2 trust. On a hit, gain 1.
    \item \textbf{Persuading:} Successful persuasions can increase trust up to +1.  
\end{proplist} 

\textbf{Losing trust:} This can be done by:
\begin{proplist}
    \item \textbf{failing missions:} Failing a mission for someone decreases trust with them. Making missions for opponents also lowers trust.
    \item \textbf{failing at work with/for someone:} Make a test of the most relevant skill or knowledge for the work against a DL to gain trust. On a miss, lose 2 trust. On a graze, lose 1.
    \item \textbf{Intimidating someone:} Loses -1 trust on a hit or crit.
    \item \textbf{Time decay:} Variable decay. Up to the players and GM
\end{proplist} 

\begin{dnote}
    \textbf{Effect on social actions:} Consider a contact's trust when deciding the IB. Choosing IB is already subjective, so there's no reason for a rule saying you should add the value of trust to IB.
\end{dnote}


The types of contacts and what they do are described below: \\

\subsection{Merchant}

Any merchant who is not a contact charges +100\% by default. Many goods and services are only provided at a certain trust or through coercion. High trust may even allow them to take payment later and give discounts. \\

% Intimidation as negotiation follows the same price scaling as negotiating, but the merchant always loses -1 trust. Intimidation as robbery can reduce Honor if that is known.

\textbf{Example Merchants}
%what do they sell at each level?


\begin{proplist}
    \item \textbf{General:} Buys and sells common goods. \\
    \item \textbf{Food:} Buys and sells food. \\
    \item \textbf{Blacksmith:} sells and repairs weapons and armor. \\
    \item \textbf{Alchemist:} sells potions and buys ingredients. \\
    \item \textbf{Mechanic:} Sells mechanisms. \\
    \item \textbf{Tamer:} sells horses, dogs, and other trained animals.\\
    \item \textbf{Smuggler:} Gets illegal items. Smuggles things into places.\\
    \item \textbf{Fence:} buys rare objects for a fair price.\\
\end{proplist}

\subsection{Mercenary} 
Accompanies players on missions for money. The higher the trust, the greatest the risk they are willing to take.  \\

\textbf{Example Mercenaries}
% Different strength levels - create in bestiary - how much for each?
\begin{proplist}
    \item \textbf{Knight:} An armored knight.  \\
    \item \textbf{Soldier:} A lightly armored warrior.\\
    \item \textbf{Archer:} An expert in bows.\\
    \item \textbf{Guide:} an expert in navigation or in the region.\\
    \item \textbf{Healer:} field healer.\\
\end{proplist}

\textbf{Mercenary morale:} adds a bonus to the mercenary's combat morale test (see Morale, Combat chapter). Starts at -2 and gets +1 per trust they have in the character.

\subsection{Service}
They trade a service or time of service for money. They are more expensive or may refuse to provide services if the level is too low. \\
% Define each service - how much at each level - 3 levels: beginner, journeyman, master

\begin{proplist}
    \item \textbf{Helper:} performs basic manual labour.\\
    \item \textbf{Healer:} speeds up healing injuries and wounds. Heals 8 IL extra per day on top of regular healing, costs 1 GP per day. Higher level: reattach limb. \\
    \item \textbf{Spy:} knows information about a certain place. Paid per information.\\
    \item \textbf{Banker:} lends money. Pays based on reputation.\\
    \item \textbf{Assassin:} kills people for money.\\
    \item \textbf{Thief:} Steals things for money.\\
\end{proplist}

\subsection{Influencer}
Influencer's role is to increase a character's popularity in town. \textbf{Influencer Strength:} ranges from 0 - 10.   \\

\begin{proplist}
    \item \textbf{Bard:} increases popularity. Improves effectiveness of celebration.\\
    \item \textbf{Clergyman:} increases popularity within the church. \\
    \item \textbf{Noble:} grants access to places, missions. \\
    \item \textbf{Employer:} pays for a job. \\
    \item \textbf{Guildmaster:} Gives access to a group of merchants, services or mercenaries. \\
\end{proplist}


\subsection{Trainer}

Trainers have a set of skills, knowledges and abilities that they can teach and a maximum level that they can help someone reach.

High level trainers will not give away their best secrets without high trust.

A trainer has a cost per day of instruction. They will charge extra without good trust or fear.

Standard rates for training are given per level of training:

\begin{stattable}{Training costs}{@{}lccccc@{}}
Level & 1 & 2 & 3 & 4 & 5  \\
\midrule
\rowcolor{fallowtint}
Price & 4 CP & 8 CP & 2 SP & 4 SP & 1 GP  \\[3pt]
\end{stattable}



\section{Missions}

A mission is what sets a character's Renown. A mission has a base renown, which is a property of the mission and does not change. The effective renown can change and is what is actually used by characters. 

\textbf{Mission requirements: }Missions have a Renown and Honor requirement. They may also require trust of a quest giver or a title to be given. 

\textbf{Mission record:} The method and quality of completion influence how much Renown the mission gives. The mission with the highest base Renown that a player is recognized for represents their Renown. Failing a mission changes the effective Renown, not the base. That means that a big failure can eclipse a former smaller success.

\textbf{Reputation decay:} A mission completed long ago has its effective Renown lowered. The rate of decay is subjective and can be removed by staging recurrent celebrations.

\textbf{Mission rewards:} Missions give rewards in the shape of karma and often money. They also often improve trust with the quest giver.

\textbf{Usurping missions:} The team can claim to have completed a mission by celebrating. The actual people who completed the mission will not like this. 


\textbf{Mission skeleton example: }

Requirement: Trust 1 with Noble, Renown: 2, Honor > -6. \\
Reputation: Renown:  3. \\
Increases trust with noble by 1 (max 2). reduces trust with merchant by -1. \\
Yield: 3 karma per participant and 20 GP in total. Yields +1 karma, +10 GP for extra objective. -1 Honor if violence happens.\\
No karma for characters with a skill > 5.\\
Failure: Loss of trust with mission giver -1. Mission renown -1. 

\subsection{Deeds}

Deeds are what define a character's honor. Deeds are crimes, acts of charity, heroism, cruelty, extortion, and other morally important marks that stay in one's reputation. They do not need to be of high consequence and difficulty, like renown. 

\textbf{Deed priority:} Deeds have an importance value. The most important deed is the measure of honor of a character.

\textbf{Gaining deeds:} doing morally and ethically relevant things can create a deed. The deed's importance is the magnitude and its honor is how morally good or bad it is.

\textbf{Losing deeds:} To remove a deed, one needs to clear their name. That means challenging a slanderer, or celebrating to disassociate name from deed. Doing something good or bad afterwards does not remove an old deed.

\textbf{Deed decay:} A deed's importance decays with time, whether it was good or bad. The decay stays within 3 points of the original importance. Celebrating to remember it will increase its importance back.


\section{Downtime}

\newcommand{\downtime}[6]{
\statrule
{\sffamily\bfseries\large Activity: #1} \\
\statlabel{Days:} #2 \\
\statlabel{Cost:} #3 \\
\statlabel{Requirement:} #4 \\
\statlabel{Test} #5 \\
\statlabel{Yield:} #6 \\
}

Players can do many different downtime activities. Each one costs a number of days and possibly money.

Only one thing can be done per day, but a single thing can fit more than one category.


\downtime{Training}{10 x N}{Trainer rate x days }{Trainer}{ none }{ N XP for skill or knowledges. }

\downtime{Training without trainer}{10 x N}{ variable }{ - }{ 2x talent vs 2x training level }{ 
    miss = 0 XP \\ 
    graze = $N/4$ XP \\
    hit $N/2$ XP \\
    crit N XP.
}

\downtime{Job}{ N days }{ - }{Employer}{ arbitrary skill or knowledge vs DL }{ 
    daily rate x N,  \\
    Gain +1 trust with employer on a crit.  \\
    If skill test DL is equal or higher than 2x training level, this counts as training alone. \\
    A miss loses 2 trust. A graze loses 1 trust.
}

\downtime{Celebrate}{ lasts 10 days }{ Living standard x 10. Can be more expensive for larger regions. }{ Populated town and a deed or mission. Only one deed or mission can be celebrated at a time, unless the deed was part of the mission.}{ Persuasion or deception vs DL. DL increases for lack of evidence and lack of previous renown. DL decreases by 1 by increasing duration by 10 days, with a maximum of 30 days. It is possible to extend the celebration after rolling for it.}{
    Improves popularity in town while celebrating to 2/4/7/10, if that is an improvement. \\
    Gives a bonus to the renown of a mission by -1/0/0/+1 depending on the test result. The bonus does not affect the base value.  \\
    Uses deception to usurp someone else's mission's renown and Honor on a hit or crit. Loses trust with the usurped and, on misses or grazes, reduces own Honor by -1/-2. \\
    If persuasion or deception DL is equal or higher than 2x training level, this counts as training alone for that skill.
}

\downtime{Healing}{N days}{rate per day}{Healer}{ }{ 
    Healer's heal rate x N in IL and wound. Payment is separate per IL and wound. The character still heals their normal daily rate.
}

\textbf{Healing alone:} Use the rules of long camping.

\subsection{ Cost of Living}

Costs of living include food, water, housing and security. 

The cost of living changes from town to town. This table is an average.

\begin{stattable}{Cost of Living}{@{}lccc@{}}
Level & Cost & Exposure & Renown   \\
\midrule
Miserable & 1CP/day & -3 & zero  \\
\rowcolor{fallowtint}
Poor & 5CP/day & -1 & 2 \\
Medium & 1SP/day & 0 & 4 \\
\rowcolor{fallowtint}
Rich & 3SP/day & 0 & 7 \\
Lavish & 1GP/day & 0 & 10 \\
\end{stattable}

A character has to live at their renown level or higher. Living lower for a month will limit the renown to the value on the table.

Higher standard of living makes for better celebrations. The recognized renown of a mission will not exceed that of the cost of living the celebration happened in.

\end{multicols}
